\documentclass[11pt,a4paper]{article}
\usepackage[T1]{fontenc}
\usepackage[utf8]{inputenc}
\usepackage{libertinus}
\usepackage[margin=26mm,headheight=14pt]{geometry}
\usepackage{amsmath,amssymb,amsthm,mathtools}
\usepackage{microtype,booktabs,array,longtable}
\usepackage{xcolor,fancyhdr,xurl,enumitem}
\usepackage[colorlinks=true,linkcolor=blue!45!black,citecolor=blue!45!black,urlcolor=blue!45!black]{hyperref}
\usepackage[nameinlink,noabbrev]{cleveref}
\setlength{\emergencystretch}{2em}
\setlength{\parskip}{0.22em}
\numberwithin{equation}{section}
\setcounter{tocdepth}{1}
\newtheorem{theorem}{Theorem}[section]
\newtheorem{lemma}[theorem]{Lemma}
\newtheorem{proposition}[theorem]{Proposition}
\newtheorem{corollary}[theorem]{Corollary}
\theoremstyle{definition}
\newtheorem{definition}[theorem]{Definition}
\newtheorem{question}{Research question}
\theoremstyle{remark}
\newtheorem{remark}[theorem]{Remark}
% ed. (2026-09-29): unnumbered environment for editorial notes.
\newtheorem*{ednote}{Editorial note (ProveIt, 2026-09-29)}
\DeclareMathOperator{\Var}{Var}
\DeclareMathOperator{\Law}{Law}
\DeclareMathOperator{\supp}{supp}
\DeclareMathOperator{\TV}{TV}
\newcommand{\R}{\mathbb R}
\newcommand{\E}{\mathbb E}
\newcommand{\Prob}{\mathbb P}
\newcommand{\N}{\mathbb N}
\newcommand{\cA}{\mathcal A}
\newcommand{\cX}{\mathcal X}
\newcommand{\norm}[1]{\left\lVert#1\right\rVert}
\newcommand{\abs}[1]{\left|#1\right|}
\newcommand{\up}{\operatorname{up}}
\newcommand{\sinc}{\operatorname{sinc}}
\newcommand{\ind}{\mathbf 1}
\pagestyle{fancy}
\fancyhf{}
\fancyhead[L]{\small Flat boundaries and uniform-factor recovery}
\fancyhead[R]{\small\thepage}
\renewcommand{\headrulewidth}{0.35pt}
\hypersetup{pdftitle={Flat Boundaries Preserve Information: Uniform-Factor Recovery without Gaussian Smoothing},pdfauthor={Research report prepared for Vladimir Reshetnikov},pdfsubject={Fabius--Rvachev density, Hellinger asymptotics, minimax estimation}}

\begin{document}
% ed. (2026-09-29): no page anchor on the title page (the build reported a
% duplicate destination page.1).
\hypersetup{pageanchor=false}
\begin{titlepage}
\centering
\vspace*{15mm}
{\Large\scshape ProveIt research continuation\par}
\vspace{13mm}
{\huge\bfseries Flat Boundaries Preserve Information\par}
\vspace{5mm}
{\Large Sharp recovery of uniform convolution factors\par
without Gaussian smoothing\par}
\vspace{12mm}
{\large Research report prepared for Vladimir Reshetnikov\par}
\vspace{2mm}
{29 September 2026\par}
\vspace{13mm}
\begin{minipage}{0.92\textwidth}
\begin{abstract}
A preceding ProveIt report established total-variation stability for finite
uniform-factor deconvolution but left statistical lower bounds unresolved.
A companion report obtained sharp rates under a strictly positive Gaussian
variance floor and explicitly asked whether that floor could be removed for
the Fabius--Rvachev density. We supply the missing analytic argument.
For every infinite summable convolution of positive-width centered uniforms,
and every pair of perturbation laws having all absolute moments finite, a
first unmatched moment of order $r$ gives an exact squared Hellinger
asymptotic of order $h^{2r}$, with constant
$\Delta_r^2\int (f^{(r)})^2/f\,/\bigl(4(r!)^2\bigr)$.
Finite uniform prefixes, conditional-score contraction, and convergence of
higher-order information provide a proof that controls the moving support
rather than ignoring it.

Consequently the sharp scale-recovery rate is $n^{-1/(4m)}$ with known
Gaussian variance, including variance zero. With unknown variance in
$[0,V]$, the scale and variance rates are $n^{-1/(4m+4)}$ and
$n^{-1/(2m+2)}$; matching hard pairs approach the zero-noise boundary and
one member has exactly zero Gaussian variance. We also prove a finite-prefix
boundary trichotomy, including an explicit logarithmic critical case, and
exhibit identifiable perturbations whose Hellinger separation is smaller than
every power of their amplitude. The article includes full conventional
proofs, reproducible checks, and a research agenda. No proof-assistant
verification or worldwide publication-priority claim is made.
\end{abstract}
\end{minipage}
\vfill
{\small Repository snapshot:\par
\texttt{9b24a3a8d545af9624f6ac455f5b548be62818b6}\par}
\end{titlepage}
% ed. (2026-09-29): page anchors resume after the title page.
\hypersetup{pageanchor=true}
\tableofcontents
\newpage

\section{Research target, provenance, and contribution boundary}
\label{sec:scope}

The question is not whether moment cancellation occurs. It does. The question
is whether cancellation survives division by a density that vanishes to all
orders at the endpoints of its support. This is the step needed to turn an
algebraic inverse modulus into a statistically sharp lower bound.

\subsection{The specific gap being continued}
At the pinned repository snapshot, the report \emph{Sharp Stability Strata
for Finite Fabius--Rvachev Deconvolution} proves local total-variation inverse
exponents, exact density expansions, and a sampling upper bound
\cite{repo-strata}. Its source, lines 1190--1234, explicitly declines to call
that upper bound statistically minimax and identifies Hellinger or relative
entropy control as the missing ingredient. The global scale upper bound has
exponent $1/(4m)$ when the capacity is $m$.

A separate, earlier-today companion manuscript, \emph{Gaussian Confounding
and Sharp Recovery of Uniform Convolution Factors}, proves corresponding
statistical rates with a strictly positive Gaussian variance floor
\cite{companion}. Its Research Question 1, page 17, asks for the unsmoothed
up-law lower bound and specifically for control of
$\int (g^{(k)})^2/g$ and the changing support. That manuscript was read from
the user's Library, not asserted to be a file in the pinned repository.

The present article resolves this question for the up density and, more
generally, for every infinite summable positive-width uniform convolution.
It also proves that the unknown-variance lower-bound obstruction persists
arbitrarily close to zero Gaussian variance. This last assertion is stronger
than a global lower bound obtained merely by restricting to a positive-noise
submodel. It is not a claim that every mixed collision stratum is classified.

\subsection{What is classical and what is developed here}
The Fabius--Rvachev infinite-product construction is classical; see Arias de
Reyna \cite{arias}. Higher-order Fisher information is also classical.
Bobkov's treatment \cite{bobkov} discusses its origin, convolution
monotonicity, and sufficient conditions for finiteness. In particular,
all-order information finiteness for an infinite uniform convolution is not
presented here as a newly invented functional or an isolated breakthrough.
We give the elementary spline proof needed by our argument.

Moment matching and singular moment inversion have substantial statistical
precedents, including the mixture-estimation work of Heinrich--Kahn and
Wu--Yang \cite{heinrich-kahn,wu-yang}. Our parameters describe a
\emph{convolution of independent factors}, not a mixture with free weights.
Those papers provide context, not substitute proofs for this model. The
shifted-power-sum inequality in \cref{sec:algebra} is credited to the companion
manuscript and reproduced to make the statistical conclusions self-contained.

The main development here is the exact all-order Hellinger theorem at a flat
moving boundary, its finite-prefix boundary trichotomy, and the resulting
zero-noise statistical lower bounds. The proofs are mathematical arguments
in this article, not kernel-checked results. The literature comparison is
focused; it establishes a precise continuation of the identified reports,
not exhaustive worldwide novelty.

\section{Definitions and main results}
\label{sec:main}

Let $(U_j)_{j\ge1}$ be independent uniforms on $[-1,1]$. Fix
\begin{equation}
 b_j>0,\qquad S=\sum_{j\ge1}b_j<\infty,\qquad
 B=\sum_{j\ge1}b_jU_j.
 \label{eq:background}
\end{equation}
All infinite sums here converge absolutely for every outcome. Write $f$ for
the density of $B$, and $f_N$ for that of
$B_N=\sum_{j=1}^N b_jU_j$. We will prove the smoothness and positivity facts
used below. The Fabius--Rvachev specialization is
\begin{equation}
 b_j=2^{-j},\qquad S=1,\qquad f=\up.
 \label{eq:up}
\end{equation}
With the characteristic-function convention $\E e^{itB}$, it has product
$\prod_{j\ge1}\sinc(2^{-j}t)$.

Our Hellinger normalization, throughout, is
\begin{equation}
 H^2(p,q)=\int_\R(\sqrt p-\sqrt q)^2\,dx,
 \qquad \rho(p,q)=\int_\R\sqrt{pq}\,dx=1-\tfrac12H^2(p,q).
 \label{eq:hellinger}
\end{equation}
Thus $0\le H^2\le2$ and $\TV(p,q)=\frac12\int|p-q|\le H(p,q)$.
For a probability law $\nu$, let $\nu_h$ denote the law of $hZ$ when
$Z\sim\nu$, with $h>0$. A law is \emph{all-moment} if
$\int |z|^k\,d\nu(z)<\infty$ for every integer $k\ge1$. An exponential
moment is not required.

\begin{theorem}[Exact Hellinger moment asymptotic]
\label{thm:main-hellinger}
Let $f$ be defined by \eqref{eq:background}. For every integer $r\ge1$,
\begin{equation}
 0<J_r(f):=\int_{f>0}\frac{(f^{(r)}(x))^2}{f(x)}\,dx<\infty.
 \label{eq:J}
\end{equation}
Suppose $\nu,\eta$ are all-moment laws whose moments of orders
$0,\ldots,r-1$ agree, and set
$\Delta_r=\int z^r\,d\nu-\int z^r\,d\eta$. Then
\begin{equation}
 H^2(f*\nu_h,f*\eta_h)
 =\frac{\Delta_r^2}{4(r!)^2}J_r(f)h^{2r}+o(h^{2r}).
 \label{eq:main-H}
\end{equation}
More precisely, in $L^2(\R)$,
\begin{equation}
 \frac{\sqrt{f*\nu_h}-\sqrt{f*\eta_h}}{h^r}
 \longrightarrow
 \frac{(-1)^r\Delta_r}{2r!}\frac{f^{(r)}}{\sqrt f},
 \label{eq:main-root}
\end{equation}
where the right side is zero outside $(-S,S)$. The statement includes
$\Delta_r=0$, in which case \eqref{eq:main-H} is a little-oh assertion.
\end{theorem}

For the statistical model, fix $m\ge1$ and set
\begin{equation}
 \cA_m=\{a\in[0,1]^m:a_1\le\cdots\le a_m\},\qquad
 X_{a,v}=B+\sum_{i=1}^m a_i V_i+\sqrt v\,G,
 \label{eq:model}
\end{equation}
where the $V_i$ are fresh uniforms on $[-1,1]$, $G$ is standard normal, and
all variables are independent. Let $P_{a,v}=\Law(X_{a,v})$.
Zeros pad the factor list to a fixed capacity; they are not observable
positive factors. The background law is known.

\begin{theorem}[Sharp estimation, including the zero-noise boundary]
\label{thm:minimax}
For $n$ independent observations, define minimax expected loss using
$\norm{\widehat a-a}_\infty$ for scales and $|\widehat v-v|$ for variance.
With known fixed $v_0\ge0$,
\begin{equation}
 \inf_{\widehat a}\sup_{a\in\cA_m}
 \E_{a,v_0}\norm{\widehat a-a}_\infty\asymp n^{-1/(4m)}.
 \label{eq:known-rate}
\end{equation}
With unknown $v\in[0,V]$, for fixed $V>0$,
\begin{align}
 \inf_{\widehat a}\sup_{a\in\cA_m,\,0\le v\le V}
 \E_{a,v}\norm{\widehat a-a}_\infty&\asymp n^{-1/(4m+4)},
 \label{eq:unknown-rate}\\
 \inf_{\widehat v}\sup_{a\in\cA_m,\,0\le v\le V}
 \E_{a,v}|\widehat v-v|&\asymp n^{-1/(2m+2)}.
 \label{eq:variance-rate}
\end{align}
The constants may depend on $m$, the fixed background, and $v_0$ or $V$.
The lower bounds in \eqref{eq:unknown-rate}--\eqref{eq:variance-rate} hold
already on every fixed relative neighborhood of $(a,v)=(0,0)$, using pairs
for which one variance is exactly zero. In \eqref{eq:known-rate}, hard pairs
approach $a=0$, also when $v_0=0$.
\end{theorem}

The last clause is a local \emph{minimax} statement. It does not assert that
the risk of every estimator at the single parameter $(0,0)$ has a lower
bound: the identically zero estimator has zero loss at that one point.

\begin{theorem}[Detection thresholds]
\label{thm:detection-main}
Test no positive factors, $a=0$, against $a_m\ge h$, with fixed capacity
$m$. With known $v_0\ge0$, the critical separation order is $n^{-1/4}$:
consistent testing is possible if $nh^4\to\infty$, and the optimal sum of
worst-case type-I and type-II errors tends to one if $nh^4\to0$.
With unknown $v\in[0,V]$, replace $nh^4$ by $nh^8$; the critical order is
$n^{-1/8}$. Both assertions include alternatives having exactly zero
Gaussian variance.
\end{theorem}

\section{Uniform prefixes and higher-order information}
\label{sec:information}

\subsection{Smoothness, support, and the endpoint polynomial}
Let $A_N=\sum_{j\le N}b_j$ and $s_N=S-A_N$. Repeated integration of interval
indicators shows that $f_N$ is a compactly supported spline of degree
$N-1$, is positive on $(-A_N,A_N)$, and belongs to $C^{N-2}$ for $N\ge2$.
These claims can also be read directly from the formula
\begin{equation}
 f_N(x)=\frac{1}{(N-1)!\prod_{j=1}^N2b_j}
 \sum_{E\subseteq\{1,\ldots,N\}}(-1)^{|E|}
 \left(x+A_N-2\sum_{j\in E}b_j\right)_+^{N-1}.
 \label{eq:spline}
\end{equation}
For $N=1$ the power zero means an interval step, with its value at a knot
irrelevant. The formula follows inductively from convolution with
$(2b_N)^{-1}\ind_{[-b_N,b_N]}$. Positivity follows by integrating a positive
previous density over the nonempty overlap of the two support intervals.

For later use the endpoint has the particularly simple expression
\begin{equation}
 f_N(A_N-d)=f_N(-A_N+d)=C_Nd^{N-1},\quad
 C_N=\frac{1}{(N-1)!\prod_{j=1}^N2b_j},
 \quad 0<d<2\min_{j\le N}b_j.
 \label{eq:endpoint}
\end{equation}
Indeed, the deficits $b_j-b_jU_j$ are uniform on $[0,2b_j]$; at such small
total deficit their convolution is the unconstrained simplex-volume
formula.

\begin{lemma}[Regularity of the infinite background]
\label{lem:smooth}
The law in \eqref{eq:background} has a $C^\infty$ density $f$, supported on
$[-S,S]$, positive on $(-S,S)$. For every fixed $k\ge0$,
$f_N^{(k)}\to f^{(k)}$ uniformly once $N$ is large enough for these classical
derivatives to exist. All derivatives of $f$ vanish at both endpoints.
\end{lemma}
\begin{proof}
Fix $L\ge k+2$. Convolution of $f_L$ with the law of the remaining infinite
sum gives a density with bounded continuous derivatives through order $k$.
This represents the same law for each $k$, so gives $f\in C^\infty$.
Couple the finite and infinite tails after $L$ with the same uniforms.
Their difference is bounded by $s_N$. Uniform continuity of $f_L^{(k)}$
therefore bounds $\norm{f_N^{(k)}-f^{(k)}}_\infty$ by its modulus of
continuity at $s_N$, which tends to zero.

For $|x|<S$, choose $N$ so large that $|x|+s_N<A_N$. Every shift of $x$ by
an element of the tail support is then inside the positive interval of
$f_N$. Thus $f(x)=\E f_N(x-(B-B_N))>0$. Support inclusion follows directly
from $|B|\le S$; positivity proves equality of the support. Smoothness and
identical vanishing outside the support force every endpoint derivative to
vanish.
\end{proof}

\subsection{A conditional-score inequality}
For integer $r\ge1$ and real $q\ge2$, put
\begin{equation}
 J_{r,q}(g)=\int_{g>0}\left|\frac{g^{(r)}}g\right|^qg,
 \qquad J_r(g)=J_{r,2}(g).
 \label{eq:Jrq}
\end{equation}
We apply this only to densities whose indicated weak derivative is an
integrable function and is zero almost everywhere on their zero set.
All sufficiently long prefixes satisfy these conditions.

\begin{lemma}[Convolution contraction]
\label{lem:score}
For such a density $g$ and any probability law $\mu$,
\begin{equation}
 J_{r,q}(g*\mu)\le J_{r,q}(g)
 \label{eq:contraction-J}
\end{equation}
whenever the right side is finite.
\end{lemma}
\begin{proof}
Write $X\sim g$, $Y\sim\mu$, independently. Differentiating the convolution
in the weak sense gives $(g*\mu)^{(r)}=g^{(r)}*\mu$. The ratio of this
identity to $g*\mu$ says, for almost every $z$ of positive density,
\[
 \frac{(g*\mu)^{(r)}(z)}{(g*\mu)(z)}
 =\E\left[\frac{g^{(r)}(X)}{g(X)}\,\middle|\,X+Y=z\right].
\]
Conditional Jensen followed by expectation proves the claim.
\end{proof}
This is the standard score-contraction mechanism; its use here does not
claim a new general information inequality.

\begin{proposition}[All-order finiteness and convergence of prefixes]
\label{prop:Jconvergence}
For every $r\ge1$, $q\ge2$, $J_{r,q}(f)<\infty$. Moreover,
\begin{equation}
 J_r(f_N)\downarrow J_r(f)\quad\text{as }N\to\infty,
 \label{eq:Jlimit}
\end{equation}
where the sequence is considered from any index $N>2r$ onward. The limit
is strictly positive.
\end{proposition}
\begin{proof}
If $N>rq$, the endpoint formula gives
\[
 \left|\frac{f_N^{(r)}(A_N-d)}{f_N(A_N-d)}\right|^q
 f_N(A_N-d)=O(d^{N-1-rq}).
\]
This is integrable; away from the endpoints the density is bounded below
and its indicated derivatives are bounded. Thus $J_{r,q}(f_N)<\infty$.
Convolution with the infinite tail and \cref{lem:score} prove finiteness for
$f$. Convolution with one further uniform proves monotonicity of $J_r(f_N)$.

To identify its limit, fix $q>2$ and a prefix $L>rq$. For all $N\ge L$,
$J_{r,q}(f_N)\le K:=J_{r,q}(f_L)$. On every compact interval inside
$(-S,S)$ the integrands defining $J_r(f_N)$ converge uniformly to that of
$J_r(f)$, by \cref{lem:smooth}. On the complementary endpoint region
$E_\delta=\{x:|x|>S-\delta\}$, H\"older's inequality gives
\begin{equation}
 \int_{E_\delta}\frac{(f_N^{(r)})^2}{f_N}
 \le K^{2/q}\Prob(|B_N|>S-\delta)^{1-2/q}.
 \label{eq:uniform-integrability}
\end{equation}
When $s_N<\delta$, the probability is at most
$\Prob(|B|>S-2\delta)$ by the deterministic coupling bound
$|B_N-B|\le s_N$. This tends to zero as $\delta\downarrow0$, since $B$
has no endpoint atoms. The same estimate applies to $f$. Interior convergence
and these tail estimates prove \eqref{eq:Jlimit}.

If $J_r(f)=0$, then $f^{(r)}$ is zero throughout $(-S,S)$ by continuity.
It is zero outside as well, so $f$ is a polynomial of degree at most $r-1$
on $\R$. A compactly supported polynomial is zero, contradicting integral
one. Hence $J_r(f)>0$.
\end{proof}

\section{The finite-prefix Hellinger expansion}
\label{sec:finite-expansion}

The following lemma isolates the analytic issue. It deliberately requires
more uniform factors than twice the first potentially unmatched moment
order; \cref{sec:phase} will show why the threshold is sharp.

\begin{lemma}[A sufficiently long finite prefix]
\label{lem:finite-H}
Let $g$ be the density of a sum of $M$ independent positive-width centered
uniforms, with $M>2r$. If all-moment laws $\nu,\eta$ agree in moments
$0,\ldots,r-1$, then
\begin{equation}
 H^2(g*\nu_h,g*\eta_h)
 =\frac{\Delta_r^2}{4(r!)^2}J_r(g)h^{2r}+o(h^{2r}).
 \label{eq:finite-H}
\end{equation}
\end{lemma}
\begin{proof}
Write $[-A,A]$ for the support. First suppose the laws are supported on
$[-L,L]$, and let $p_h=g*\nu_h$, $q_h=g*\eta_h$.
On a fixed compact subinterval of $(-A,A)$, Taylor expansion and moment
cancellation give
\begin{equation}
 h^{-r}(p_h-q_h)\longrightarrow
 \frac{(-1)^r\Delta_r}{r!}g^{(r)}.
 \label{eq:pointwise-finite}
\end{equation}
The available smoothness is sufficient since $M>2r$ implies $g\in C^r$.
Both densities tend to $g$, giving the claimed pointwise square-root limit.

Consider a small endpoint collar and write $d=A-|x|$. On the interior part
$d\ge2hL$, all shifted arguments stay inside the polynomial endpoint collar
provided its fixed width is reduced once. There the endpoint formula and
Taylor's integral remainder yield
\begin{equation}
 p_h(x)+q_h(x)\ge c d^{M-1},\qquad
 |p_h(x)-q_h(x)|\le C h^r d^{M-1-r}.
 \label{eq:collar-bound}
\end{equation}
Constants may depend on the fixed laws and prefix. Consequently
\begin{equation}
 h^{-2r}(\sqrt{p_h}-\sqrt{q_h})^2
 \le C d^{M-1-2r}.
 \label{eq:collar-domination}
\end{equation}
The exponent is greater than $-1$. On the remaining fixed interior the same
quantity is bounded, because $g$ is bounded below and its $r$th derivative
is bounded. Dominated convergence is therefore valid away from the shrinking
boundary band. That band, including the new support outside $[-A,A]$, has
width $O(h)$ and total $p_h+q_h$ mass $O(h^M)$. Its Hellinger contribution
is at most that mass, and $h^M=o(h^{2r})$.

We give the extension to unbounded all-moment laws explicitly. Choose
$\varepsilon>0$ with $M(1-\varepsilon)>2r$ and put $L_h=h^{-\varepsilon}$.
Condition each law on $|Z|\le L_h$ to obtain $\nu^{[h]},\eta^{[h]}$.
For every prescribed $K$, the discarded mass and each of the finitely many
discarded moments through order $r$ are $O(h^K)$: apply Markov's inequality
with a sufficiently high absolute moment. The same holds for the effect of
normalization. In particular, their moment discrepancies of orders $j<r$
are $O(h^K)$, and their discrepancy at order $r$ tends to $\Delta_r$.

Replacing a law by this conditioning changes the Hellinger distance after
any convolution by $O(h^{K/2})$. To see this, couple the original variable
and its conditional version so that they differ with probability at most
the discarded mass, and use $H^2\le2\TV$. Choose $K>2r$, so the change
in the Hellinger distance is $o(h^r)$.

For the conditioned densities use the band of width $2hL_h$. Its mass is
$O((hL_h)^M)=o(h^{2r})$. On the remaining endpoint collar the term coming
from the $r$th remainder is still bounded as in
\eqref{eq:collar-domination}, with a uniform constant: the bound uses the
uniformly bounded $r$th absolute moments, not $L_h^r$. A residual moment
of order $j<r$ contributes at most
$C h^{j+K}d^{M-1-j}$ to the density difference. After division by a
square-root lower bound and by $h^r$, its $L^2$ norm is at most
\[
 C h^{j+K-r}
 \left(\int_0^{d_0}d^{M-1-2j}\,dd\right)^{1/2}=o(1).
\]
The fixed interior has the same bound without the endpoint weight. Thus the
residual moments do not affect the limiting norm. The $r$th term is handled
by dominated convergence exactly as before. Finally the triangle inequality
for $H$, and its $o(h^r)$ conditioning error, transfer the expansion to the
original laws.
\end{proof}

The proof does not suppose that $p_h/g$ is bounded everywhere. Such a
supposition would fail even for a single added uniform, because its support
extends beyond that of $g$.

\section{Passing from prefixes to the infinite convolution}
\label{sec:infinite}

\subsection{A lower bound and a matching upper bound}
\begin{proof}[Proof of \cref{thm:main-hellinger}]
Finiteness and positivity of $J_r(f)$ are \cref{prop:Jconvergence}. Smoothness,
bounded derivatives, and the $(r+1)$st absolute moments give, at every fixed
$x$,
\[
 (f*\nu_h-f*\eta_h)(x)
 =\frac{(-1)^r\Delta_r}{r!}h^rf^{(r)}(x)+O(h^{r+1}).
\]
For $|x|<S$, division by
$\sqrt{f*\nu_h(x)}+\sqrt{f*\eta_h(x)}\to2\sqrt{f(x)}$
and Fatou's lemma imply
\begin{equation}
 \liminf_{h\downarrow0}h^{-2r}H^2(f*\nu_h,f*\eta_h)
 \ge\frac{\Delta_r^2}{4(r!)^2}J_r(f).
 \label{eq:Fatou}
\end{equation}

For the upper bound fix $N>2r$ and let $\tau_N$ be the law of the infinite
tail after $N$, so $f=f_N*\tau_N$. Hellinger distance contracts under a
common convolution:
\begin{equation}
 H^2(f*\nu_h,f*\eta_h)\le H^2(f_N*\nu_h,f_N*\eta_h).
 \label{eq:Hcontract}
\end{equation}
Indeed, pointwise Cauchy--Schwarz gives
$\sqrt{(p*\tau)(q*\tau)}\ge\sqrt{pq}*\tau$, and integration shows that
affinity cannot decrease. Applying \cref{lem:finite-H} gives
\[
 \limsup_{h\downarrow0}h^{-2r}H^2(f*\nu_h,f*\eta_h)
 \le\frac{\Delta_r^2}{4(r!)^2}J_r(f_N).
\]
Now send $N$ to infinity and apply \eqref{eq:Jlimit}. This matches
\eqref{eq:Fatou} and proves the exact norm limit.

For completeness the stronger $L^2$ conclusion follows too. The scaled
square-root differences converge pointwise inside the support to the stated
function. Outside it, for any fixed positive distance $d$ from the support,
\[
 (f*\nu_h)(x)\le\norm f_\infty\Prob(|Z|\ge d/h)=O(h^K)
\]
for every $K$, and similarly for $\eta$. Thus the scaled functions converge
almost everywhere on $\R$. Their $L^2$ norms are bounded and converge to
the norm of the pointwise limit. Boundedness and almost-everywhere convergence
imply weak $L^2$ convergence: first test against bounded functions of compact
support, using uniform integrability in $L^1$ on that support, then approximate
an arbitrary $L^2$ test function. Weak convergence plus convergence of norms
in a Hilbert space gives strong convergence, proving \eqref{eq:main-root}.
\end{proof}

The order of limits matters. We first take $h\downarrow0$ at a fixed finite
prefix, then send the prefix length to infinity. No uniform estimate in both
variables has been silently assumed.

\subsection{What happens to the new support?}
\begin{proposition}[Leakage is smaller than every power]
\label{prop:leakage}
For any all-moment perturbation $Z$, independent of $B$, and every integer
$K\ge1$,
\begin{equation}
 \Prob(|B+hZ|>S)=O(h^K).
 \label{eq:leakage}
\end{equation}
Nevertheless, for a nonzero uniform perturbation,
$\chi^2(f*\Law(hU),f)=\infty$ for every $h>0$.
\end{proposition}
\begin{proof}
If $B>S-d$, then $B_N>A_N-d$, since the tail is at most $s_N$.
The finite-prefix simplex bound gives
$\Prob(B_N>A_N-d)\le C_N' d^N$ for small $d$, and increasing the constant
makes this valid for every $d\ge0$. Reflection gives the other endpoint.
Conditioning on $Z$ therefore bounds the leakage by
$C_N' h^N\E|Z|^N$. Set $N=K$.

For an added uniform the convolution is positive on $(-S-h,S+h)$, including
sets outside $[-S,S]$ of positive measure, while $f$ is zero there. It is not
absolutely continuous relative to the law with density $f$, so the indicated
chi-square divergence is infinite.
\end{proof}
Small leakage is not zero leakage. The exact Hellinger theorem is therefore
not a disguised finite-chi-square argument.

\subsection{An identifiable, beyond-all-powers consequence}
\begin{corollary}[All moments may hide every algebraic scale]
\label{cor:allmoments}
If distinct all-moment laws $\nu,\eta$ have the same moments of every
nonnegative integer order, then, for every $K\ge1$,
\begin{equation}
 H(f*\nu_h,f*\eta_h)=o(h^K).
 \label{eq:superpoly}
\end{equation}
The two convolved laws are nonetheless distinct for every $h>0$.
\end{corollary}
\begin{proof}
Apply \cref{thm:main-hellinger} with $r=K$ and $\Delta_r=0$.
For distinctness, the characteristic function of the compactly supported
background extends to an entire function and equals one at zero. Its real
zeros are therefore isolated. Equality of the convolved characteristic
functions would force those of $\nu_h,\eta_h$ to agree off that discrete
zero set, and then everywhere by continuity. Uniqueness of characteristic
functions would give $\nu=\eta$, a contradiction.
\end{proof}
An explicit standard moment-indeterminate construction uses the lognormal
density
\[
 \ell(z)=\frac{e^{-(\log z)^2/2}}{z\sqrt{2\pi}},\quad z>0,
 \qquad \ell_\theta(z)=\ell(z)\bigl(1+\theta\sin(2\pi\log z)\bigr),
 \quad 0<|\theta|<1.
\]
Both are positive probability densities with all moments. For $Y$ standard
normal and every integer $k\ge0$,
\[
 \E[e^{kY}\sin(2\pi Y)]
 =\operatorname{Im} e^{(k+2\pi i)^2/2}
 =e^{k^2/2-2\pi^2}\sin(2\pi k)=0.
\]
This verifies the construction directly; moment-indeterminate lognormal
families are classical, with further context in \cite{kleiber}.
At any amplitude $h_n=n^{-\alpha}$, $\alpha>0$, the two $n$-sample laws
have total variation tending to zero, by \eqref{eq:superpoly} and
$H^2(P^n,Q^n)\le nH^2(P,Q)$. This does not contradict the finite-uniform-factor
rates: these lognormal perturbations are outside that finite-dimensional
model.

\section{A sharp finite-boundary phase diagram}
\label{sec:phase}

An infinite uniform convolution has arbitrarily long finite prefixes.
A finite convolution does not. The resulting distinction can be quantified
exactly, including the borderline logarithm.

Let $g_M$ be the density of $\sum_{i=1}^M b_iU_i$, now with $M$ fixed, and
let $A=\sum_{i=1}^M b_i$ and $C_M$ be as in \eqref{eq:endpoint}. Suppose
$\nu,\eta$ are symmetric all-moment laws whose first unmatched moment has
even order $r\ge2$, so $\Delta_r\ne0$. Define
\begin{equation}
 \Psi_{\nu,M}(t)=\int (z-t)_+^{M-1}\,d\nu(z),\qquad
 D_M(t)=\sqrt{\Psi_{\nu,M}(t)}-\sqrt{\Psi_{\eta,M}(t)}.
 \label{eq:profiles}
\end{equation}
For $M=1$, $(z-t)_+^0$ means $\ind_{\{z>t\}}$.

\begin{theorem}[Finite-uniform boundary trichotomy]
\label{thm:phase}
Under these hypotheses, as $h\downarrow0$:
\begin{align}
 M<2r:\quad
 H^2(g_M*\nu_h,g_M*\eta_h)
 &\sim 2C_M A_M h^M,
 \label{eq:phase-sub}\\
 M=2r:\quad
 H^2(g_M*\nu_h,g_M*\eta_h)
 &\sim \frac{C_M\Delta_r^2}{2}
 \binom{2r-1}{r}^{\!2}h^{2r}\log(1/h),
 \label{eq:phase-critical}\\
 M>2r:\quad
 H^2(g_M*\nu_h,g_M*\eta_h)
 &\sim\frac{\Delta_r^2}{4(r!)^2}J_r(g_M)h^{2r}.
 \label{eq:phase-regular}
\end{align}
Here $A_M=\int_\R D_M(t)^2\,dt\in(0,\infty)$ when $M<2r$.
In that case the leading constant depends on the full endpoint
profiles, not merely on $\Delta_r$.
\end{theorem}
\begin{proof}
The third case is \cref{lem:finite-H}. We prove the other two and justify
that no interior spline knot contributes at leading order.

\paragraph{Endpoint reduction.}
Choose a fixed collar width $d_0>0$ small enough that the two endpoint
collars are disjoint and \eqref{eq:endpoint} holds on a collar twice as
wide. At the right edge write $x=A+ht$. Replacing $g_M(y)$ by
$C_M(A-y)_+^{M-1}$ gives exactly
\[
 (g_M*\nu_h)(A+ht)=C_Mh^{M-1}\Psi_{\nu,M}(t)
\]
when the shifted argument is in that polynomial collar. The error for
$|x-A|\le d_0$ only involves $|Z|\ge d_0/h$, after reducing $d_0$ if
necessary. All absolute moments imply that its $L^1$ norm on the collar
is $O(h^K)$ for every prescribed $K$; the extra polynomial factor of degree
$M-1$ is absorbed by taking a higher moment. The same statement holds for
$\eta$. Since $\norm{\sqrt p-\sqrt q}_2^2\le\norm{p-q}_1$ for
nonnegative functions, these errors are negligible to any prescribed
algebraic order in the Hellinger calculation.

Thus the right collar contributes, to the required accuracy,
\begin{equation}
 C_Mh^M\int_{-d_0/h}^{d_0/h}D_M(t)^2\,dt.
 \label{eq:collar-profile}
\end{equation}
The left collar has the same value by symmetry. Outside both the support
and the collars, the total density mass is smaller than every power of
$h$, again by a high-moment tail estimate.

\paragraph{Asymptotics of the profiles.}
Put $k=M-1$ and let $u\to\infty$. Replacing $(u+Z)_+^k$ by
$(u+Z)^k$ creates an error $O(u^{-K})$ for every prescribed $K$, by the
all-moment assumption. If $k\ge r$, binomial expansion and matching of
moments below $r$ give
\begin{equation}
 D_M(-u)\sim\frac{\Delta_r}{2}\binom{M-1}{r}
 u^{(M-1)/2-r}.
 \label{eq:profile-tail}
\end{equation}
If $k<r$, all polynomial terms cancel and $D_M(-u)$ decays faster than
any inverse power, after division by its denominator of order $u^{k/2}$.
For $t\to+\infty$, each profile itself decays faster than any inverse
power, since
$\E (Z-t)_+^k\le t^{-K}\E|Z|^{k+K}$ for $t>0$.

Consequently $D_M^2$ is integrable when $M<2r$. Its integral is positive:
if the profiles agreed almost everywhere, their $M$th distributional
derivatives would agree, whereas
\[
 \frac{d^M}{dt^M}\Psi_{\nu,M}
 =(-1)^M(M-1)!\,\nu
\]
as distributions. This would imply $\nu=\eta$, contrary to
$\Delta_r\ne0$. This proves the constant in \eqref{eq:phase-sub}.
When $M=2r$, \eqref{eq:profile-tail} instead gives
\[
 D_M(-u)^2\sim\frac{\Delta_r^2}{4}
 \binom{2r-1}{r}^{\!2}\frac1u.
\]
Integrating to $d_0/h$ in \eqref{eq:collar-profile}, and adding the two
endpoints, gives \eqref{eq:phase-critical}.

\paragraph{The interior.}
On a compact subinterval of $(-A,A)$ the two convolved densities are bounded
below for small $h$. If $M\ge r+1$, the spline belongs to
$W^{r,\infty}$; Taylor's integral bound and moment cancellation give a
uniform density difference $O(h^r)$ there. Its Hellinger contribution is
$O(h^{2r})$, which is negligible in the first two regimes.

If $2\le M\le r$, use the finite truncated-power representation
\eqref{eq:spline}. Since all moments through degree $M-1$ match, the
difference produced by a single knot $\tau$ is $h^{M-1}$ times a profile
in $(x-\tau)/h$ that decays faster than any inverse power in both
directions. Its squared $L^2$ integral is $O(h^{2M-1})$.
There are finitely many knots, even when subset sums coincide; the squared
norm of their sum has the same order up to a fixed constant. Division by
the interior density lower bound preserves this estimate. As
$2M-1>M$ for $M\ge2$, it is negligible compared with $h^M$.
For $M=1$ there are no interior knots and the density is constant in the
interior; only high-moment tails can affect that fixed interior region.
This completes the proof.
\end{proof}

\begin{corollary}[Simple testing scales for a finite prefix]
\label{cor:finite-testing}
For the simple pair in \cref{thm:phase}, the critical amplitude orders are
\begin{equation}
 \begin{cases}
 n^{-1/M},&M<2r,\\
 (n\log n)^{-1/(2r)},&M=2r,\\
 n^{-1/(2r)},&M>2r.
 \end{cases}
 \label{eq:finite-testing}
\end{equation}
Precisely, the product affinity tends to one when $nH^2\to0$, to zero
when $nH^2\to\infty$, and to $e^{-\lambda/2}$ when $nH^2\to\lambda$.
\end{corollary}
\begin{proof}
Independence gives
$\rho(P^n,Q^n)=(1-H^2(P,Q)/2)^n$. Insert the three asymptotics.
The inequalities $\TV\le H$ and $1-\TV\le\rho$ identify the
indistinguishable and distinguishable regimes.
\end{proof}
These are simple-pair testing scales. They are not a blanket minimax theorem
for all finite-prefix parameter models. In particular, boundary information
can improve a specific testing problem even though moment fitting alone
does not use that information.

\section{Two explicit boundary transitions}
\label{sec:transitions}

First compare $\nu=\Law(U)$ with $\eta=\delta_0$. The first unmatched
moment is $r=2$, with $\Delta_2=1/3$. For $M$ background uniforms all of
half-length one,
\begin{equation}
 C_M=\frac1{2^M(M-1)!}.
 \label{eq:equal-C}
\end{equation}
At $M=1$, direct integration of the two linear edge ramps gives, for
$0<h<1$,
\begin{equation}
 H^2(g_1*\Law(hU),g_1)
 =\frac{2(\sqrt2-1)}3h.
 \label{eq:M1}
\end{equation}
At the critical number $M=4$,
\begin{equation}
 H^2(g_4*\Law(hU),g_4)
 \sim\frac1{192}h^4\log(1/h).
 \label{eq:M4}
\end{equation}
For every $M>4$, the asymptotic is $J_2(g_M)h^4/144$.
For the infinite up density, it is likewise
\begin{equation}
 H^2(\up*\Law(hU),\up)\sim\frac{J_2(\up)}{144}h^4.
 \label{eq:up-single}
\end{equation}
Thus a nonzero added support does not force an $h$-scale Hellinger
contribution: the order of boundary vanishing controls its importance.

Second compare a uniform perturbation with a variance-matched Gaussian:
$\nu=\Law(U)$ and $\eta=N(0,1/3)$. The first unmatched moment is $r=4$,
\begin{equation}
 \Delta_4=\frac15-\frac13=-\frac2{15}.
 \label{eq:delta4-single}
\end{equation}
The critical finite prefix now has $M=8$ factors, and
\begin{equation}
 H^2(g_8*\Law(hU),g_8*N(0,h^2/3))
 \sim\frac7{829440}h^8\log(1/h).
 \label{eq:M8}
\end{equation}
The infinite up-law counterpart is
\begin{equation}
 H^2(\up*\Law(hU),\up*N(0,h^2/3))
 \sim\frac{J_4(\up)}{129600}h^8.
 \label{eq:up-confounded}
\end{equation}
These constants use the normalization \eqref{eq:hellinger}. Halving the
definition of squared Hellinger distance would halve every displayed
Hellinger constant, but would not change any rate exponent.

\section{The finite algebra needed for recovery}
\label{sec:algebra}

\subsection{Cumulants of a uniform factor}
Let $p_k(x)=\sum_{i=1}^m x_i^k$, with $x_i=a_i^2$. Cumulants are
coefficients in
$\log\E e^{tX}=\sum_{j\ge1}\kappa_j(X)t^j/j!$ near zero.
For $U$ uniform on $[-1,1]$,
\begin{equation}
 c_k:=\kappa_{2k}(U)
 =\frac{2^{2k}B_{2k}}{2k}
 =\frac{(-1)^{k+1}(2k)!\zeta(2k)}{k\pi^{2k}}\ne0.
 \label{eq:cumulants}
\end{equation}
For example, $c_1=1/3$, $c_2=-2/15$, $c_3=16/63$.
One verification is to take logarithms in
$\sinh t/t=\prod_{j\ge1}(1+t^2/(\pi^2j^2))$ for $|t|<\pi$
and compare coefficients. Cumulant additivity and scaling give
\begin{align}
 \kappa_2(X_{a,v})-\kappa_2(B)&=v+\tfrac13p_1(x),
 \label{eq:kappa2}\\
 \kappa_{2k}(X_{a,v})-\kappa_{2k}(B)&=c_kp_k(x),\quad k\ge2.
 \label{eq:kappahigher}
\end{align}
Known variance reveals $p_1,\ldots,p_m$ through cumulant order $2m$.
Unknown variance conceals $p_1$ and reveals the consecutive window
$p_2,\ldots,p_{m+1}$ through order $2m+2$.

For the background \eqref{eq:up}, the explicit formulas are
\begin{equation}
 \kappa_{2k}(B)=\frac{c_k}{2^{2k}-1},\qquad
 \Var(B)=\frac19,\qquad \kappa_4(B)=-\frac2{225}.
 \label{eq:up-cumulants}
\end{equation}
More generally replace $(2^{2k}-1)^{-1}$ by $\sum_jb_j^{2k}$.
Absolute summability justifies differentiation of the logarithmic moment
product in a neighborhood of zero.

\subsection{A credited shifted-power-sum inequality}
The following result and its counting-function proof are reproduced from
\cite[Theorem 4.1]{companion}. They are not counted as a new theorem of this
article.

\begin{lemma}[Consecutive shifted power sums]
\label{lem:moments}
For ordered $x,y\in[0,1]^m$, integer $s\ge0$, and
\[
 E_s(x,y)=\max_{s+1\le k\le s+m}|p_k(x)-p_k(y)|,
\]
one has
\begin{equation}
 \norm{x-y}_\infty\le8E_s(x,y)^{1/(m+s)},\qquad
 \norm{\sqrt x-\sqrt y}_\infty
 \le\sqrt8 E_s(x,y)^{1/(2m+2s)}.
 \label{eq:moment-inverse}
\end{equation}
\end{lemma}
\begin{proof}
Define $C(t)=\#\{i:x_i\le t\}-\#\{i:y_i\le t\}$ on $(0,1)$.
The nonzero sign of this integer-valued step function changes at most
$m-1$ times. Indeed, resolve its jumps into $m$ unit up-steps and $m$
unit down-steps. The resulting walk starts and ends at zero and has at most
$m$ excursions; combining consecutive excursions of the same sign can only
reduce the number of sign blocks. Intermediate steps at a common atom have
zero interval length and create no extra sign block of $C$.

Choose $z_1,\ldots,z_q$ between opposite-sign blocks, $q\le m-1$, and
choose the sign of $Q(t)=\pm\prod_{j=1}^q(t-z_j)$ so that $QC\ge0$
almost everywhere. The sum of the absolute monomial coefficients of $Q$
is at most $2^q$.

Put $d=\norm{x-y}_\infty>0$. Interchanging the vectors if necessary,
some $y_i-x_i=d$. On $(x_i,y_i)$, $C\ge1$, so no sign-change point
lies in that interval. On its middle half, $t\ge d/4$ and
$|t-z_j|\ge d/4$. Hence
\begin{equation}
 \int_0^1t^sQ(t)C(t)\,dt\ge\frac d2\left(\frac d4\right)^{s+q}.
 \label{eq:moment-lower}
\end{equation}
Direct integration of the counting functions gives
\[
 p_k(x)-p_k(y)=-k\int_0^1t^{k-1}C(t)\,dt.
\]
Expanding $Q$ therefore bounds the left side of
\eqref{eq:moment-lower} above by $2^q E_s/(s+1)$.
Since $d\le1$ and $q\le m-1$,
\[
 d^{m+s}\le\frac{2^{2s+3m-2}}{s+1}E_s
 \le8^{m+s}E_s.
\]
This proves the first inequality; the second follows from
$|\sqrt u-\sqrt v|\le\sqrt{|u-v|}$ for nonnegative $u,v$.
\end{proof}
Equal unit multiplicities and nonnegative nodes are essential hypotheses of
this proof. Arbitrary signed or unequal-weight moment inversion is a
different problem.

\subsection{Exact moment-preserving pairs}
The construction below, also used in \cite{companion}, supplies fixed
algebraic shapes. Our new analytic theorem will determine the Hellinger
distance between their scaled statistical realizations.

\begin{lemma}[Hard shapes]
\label{lem:hard-shapes}
For every $m\ge1$ there exist distinct positive ordered vectors $u,w$ in
$(0,1)^m$ such that
\begin{equation}
 p_1(u)=p_1(w),\ldots,p_{m-1}(u)=p_{m-1}(w),\qquad
 p_m(u)\ne p_m(w).
 \label{eq:known-shapes}
\end{equation}
There also exist such vectors with
\begin{equation}
 \begin{split}
 p_2(u)&=p_2(w),\ldots,p_m(u)=p_m(w),\\
 p_1(u)&\ne p_1(w),\qquad p_{m+1}(u)\ne p_{m+1}(w).
 \end{split}
 \label{eq:unknown-shapes}
\end{equation}
Empty strings of equalities, when $m=1$, impose no constraint.
\end{lemma}
\begin{proof}
Start at distinct $0<u_1<\cdots<u_m<1$, let
$P(z)=\prod_i(z-u_i)$, and define the tangent
\[
 \beta_i=\frac1{u_i^sP'(u_i)},\quad s\in\{0,1\}.
\]
Lagrange interpolation of monomials of degrees at most $m-1$ gives
\[
 \sum_i\frac{u_i^j}{P'(u_i)}=
 \begin{cases}0,&0\le j<m-1,\\1,&j=m-1.\end{cases}
\]
Thus the differentials of $p_{s+1},\ldots,p_{s+m-1}$ annihilate $\beta$,
whereas $Dp_{s+m}(u)\beta=s+m$. The constraint differentials have rank
$m-1$, by the rectangular Vandermonde matrix after scaling nonzero rows
and columns. The implicit function theorem gives a local one-dimensional
curve preserving those moments, with tangent $\beta$; another sufficiently
nearby point gives the required change in the final moment. The coordinates
stay distinct, positive, and less than one.

For $s=1$, partial fractions of $1/P$ evaluated at zero also give
\[
 Dp_1(u)\beta=\sum_i\frac1{u_iP'(u_i)}
 =\frac{(-1)^{m+1}}{\prod_i u_i}\ne0.
\]
Choose the second point close enough that both this moment and $p_{m+1}$
change. This proves the additional assertion in \eqref{eq:unknown-shapes}.
\end{proof}

\section{A feasible estimator and the upper bounds}
\label{sec:upper}

\begin{proposition}[Finite-grid moment fitting]
\label{prop:estimator}
There are measurable estimators attaining the upper bounds in
\cref{thm:minimax}, uniformly over the stated parameter sets.
\end{proposition}
\begin{proof}
For each required raw moment order $j$ form the sample average of $X^j$.
Clip it to a fixed interval containing every true raw moment of that order
in the model. All absolute moments of $X_{a,v}$ of any fixed order are
uniformly bounded: the background and uniform factors are bounded, and the
Gaussian variance is bounded. Consequently the clipped raw-moment estimates
have uniformly bounded mean squared error $C_j/n$.

Cumulants through any fixed order are polynomials in raw moments. On the
fixed clipping box these polynomials are Lipschitz. Hence the resulting
cumulant estimates satisfy
\begin{equation}
 \E|\widehat\kappa_j-\kappa_j(X_{a,v})|^2\le C_j'/n.
 \label{eq:kappa-MSE}
\end{equation}
With known variance, use \eqref{eq:kappa2}--\eqref{eq:kappahigher} to
estimate $p_1,\ldots,p_m$. With unknown variance, estimate
$p_2,\ldots,p_{m+1}$. Write $s=0$ or $s=1$ respectively, and
$\widehat p_k$ for these estimates. Their maximum error $D_n$ has
$\E D_n\le Cn^{-1/2}$.

On the finite ordered grid
\[
 \cX_{m,n}=\{x\in\{0,1/n,\ldots,1\}^m:x_1\le\cdots\le x_m\},
\]
minimize $\max_{s+1\le k\le s+m}|p_k(x)-\widehat p_k|$,
breaking ties lexicographically. The true squared vector has an ordered
rounded grid point within $1/n$ coordinatewise; its power-sum discrepancy
is at most $C/n$. Thus the chosen $\widehat x$ satisfies
\begin{equation}
 E_s(\widehat x,x)\le2D_n+C/n.
 \label{eq:fit-error}
\end{equation}
Set $\widehat a_i=\sqrt{\widehat x_i}$. Apply \cref{lem:moments}, then
Jensen's inequality for the concave positive powers involved, to obtain
\begin{align}
 \E\norm{\widehat x-x}_\infty&\le Cn^{-1/(2m+2s)},
 \label{eq:x-rate}\\
 \E\norm{\widehat a-a}_\infty&\le Cn^{-1/(4m+4s)}.
 \label{eq:a-upper}
\end{align}

In the unknown-variance case define
\[
 \widehat v=\Pi_{[0,V]}
 \left(\widehat\kappa_2-\kappa_2(B)-\tfrac13p_1(\widehat x)\right),
\]
where $\Pi$ is clipping to the interval. Since clipping cannot increase
distance to the true $v$,
\[
 \E|\widehat v-v|
 \le Cn^{-1/2}+\frac m3\E\norm{\widehat x-x}_\infty
 \le C'n^{-1/(2m+2)}.
\]
All minimizations were over finite fixed sets, so measurability is explicit.
\end{proof}
The grid has $\binom{n+m}{m}$ points. This is a proof of an attainable
statistical rate for fixed $m$, not a proposed scalable implementation.
Exact data also determine the parameters uniquely by \cref{lem:moments}
and the variance identity; no density inversion or use of empirical total
variation is needed.

\section{Lower bounds at vanishing factors and zero variance}
\label{sec:lower}

\subsection{Known variance}
Choose $u,w$ from \eqref{eq:known-shapes}. Let
\[
 Z_u=\sum_i\sqrt{u_i}V_i,\qquad Z_w=\sum_i\sqrt{w_i}V_i,
 \qquad a(h)=h\sqrt u,\quad a'(h)=h\sqrt w.
\]
Both perturbations are symmetric. Their cumulants through order $2m-1$
agree, and the difference in their $2m$th cumulants is
$c_m(p_m(u)-p_m(w))\ne0$. Raw moments are triangular polynomials in
cumulants with coefficient one on the highest cumulant; therefore raw
moments agree through order $2m-1$, and
\begin{equation}
 \Delta_{2m}=c_m\bigl(p_m(u)-p_m(w)\bigr).
 \label{eq:delta-known}
\end{equation}
At zero Gaussian variance, \cref{thm:main-hellinger} now gives
\begin{equation}
 H^2(P_{a(h),0},P_{a'(h),0})
 \sim\frac{\Delta_{2m}^2}{4((2m)!)^2}J_{2m}(f)h^{4m}.
 \label{eq:H-known}
\end{equation}
At any fixed known $v_0>0$, a common Gaussian convolution contracts
Hellinger distance, so the same pair has $H^2=O(h^{4m})$.
Its scale separation is a fixed positive multiple of $h$.

\subsection{Unknown variance: one member exactly unsmoothed}
Choose $u,w$ from \eqref{eq:unknown-shapes}, and define
\begin{equation}
 t_u=\max\{0,(p_1(w)-p_1(u))/3\},\qquad
 t_w=\max\{0,(p_1(u)-p_1(w))/3\}.
 \label{eq:compensation}
\end{equation}
Exactly one of these numbers is zero and the other is strictly positive.
Their purpose is the identity
$t_u+p_1(u)/3=t_w+p_1(w)/3$. Let
\[
 Z_u=\sum_i\sqrt{u_i}V_i+\sqrt{t_u}G,
 \qquad Z_w=\sum_i\sqrt{w_i}V_i+\sqrt{t_w}G.
\]
The second cumulants agree by construction, cumulants $4,6,\ldots,2m$
agree by the power-sum constraints, and odd cumulants vanish. The first
unmatched raw moment is therefore
\begin{equation}
 \Delta_{2m+2}=c_{m+1}\bigl(p_{m+1}(u)-p_{m+1}(w)\bigr)\ne0.
 \label{eq:delta-unknown}
\end{equation}
Both noise laws are all-moment, so \cref{thm:main-hellinger} applies even
though only one contains a Gaussian. With
\[
 a(h)=h\sqrt u,\quad v(h)=h^2t_u,\qquad
 a'(h)=h\sqrt w,\quad v'(h)=h^2t_w,
\]
we obtain
\begin{equation}
 H^2(P_{a(h),v(h)},P_{a'(h),v'(h)})
 \sim\frac{\Delta_{2m+2}^2}{4((2m+2)!)^2}
 J_{2m+2}(f)h^{4m+4}.
 \label{eq:H-unknown}
\end{equation}
The scale separation is comparable to $h$ and the variance separation to
$h^2$. All parameters tend to $(0,0)$, and for sufficiently small $h$ lie
in $\cA_m\times[0,V]$. One variance is exactly zero at every $h$.

\subsection{The statistical conclusion}
\begin{proof}[Completion of \cref{thm:minimax}]
For any pair of laws, tensorization yields
\begin{equation}
 H^2(P^n,Q^n)=2\left[1-\left(1-\tfrac12H^2(P,Q)\right)^n\right]
 \le nH^2(P,Q).
 \label{eq:tensorization}
\end{equation}
If two parameter values have metric separation $d$, any estimator has
maximum expected loss at least
\begin{equation}
 \frac d4\bigl(1-\TV(P^n,Q^n)\bigr).
 \label{eq:two-point}
\end{equation}
Indeed, deciding which parameter is nearer to the estimate incurs loss at
least $d/2$ whenever the resulting test errs; the sum of its two error
probabilities is at least $1-\TV$.

In the known-variance pair choose $h=\gamma n^{-1/(4m)}$, with a fixed
sufficiently small $\gamma>0$. Then \eqref{eq:H-known}, contraction if
needed, and \eqref{eq:tensorization} give product total variation at most
$1/2$ for all sufficiently large $n$. Equation \eqref{eq:two-point} gives
the lower bound $cn^{-1/(4m)}$.

For the unknown-variance pair choose
$h=\gamma n^{-1/(4m+4)}$. The same argument using
\eqref{eq:H-unknown} gives lower bounds of order $h$ for scale loss and
$h^2$ for variance loss. These are exactly
\eqref{eq:unknown-rate}--\eqref{eq:variance-rate}. The pairs eventually
lie in every fixed relative neighborhood of $(0,0)$, proving the local
minimax clause. \Cref{prop:estimator} supplied all upper bounds.
\end{proof}

\subsection{Sharp deterministic Hellinger inverse exponents}
\begin{corollary}[Inverse moduli]
\label{cor:H-inverse}
On the parameter sets in \cref{thm:minimax},
\begin{align}
 \norm{a-a'}_\infty&\le C H(P_{a,v_0},P_{a',v_0})^{1/(2m)}
 &&\text{for known }v_0,\label{eq:mod-known}\\
 \norm{a-a'}_\infty&\le C H(P_{a,v},P_{a',v'})^{1/(2m+2)},
 \label{eq:mod-unknown}\\
 |v-v'|&\le C H(P_{a,v},P_{a',v'})^{1/(m+1)}.
 \label{eq:mod-var}
\end{align}
None of the three exponents can be increased in a uniform inequality,
even locally near the corresponding all-zero factor configuration.
\end{corollary}
\begin{proof}
For densities $p,q$, Cauchy--Schwarz gives
\[
 \left|\int x^j(p-q)\,dx\right|
 \le H(p,q)\left(2\E_p|X|^{2j}+2\E_q|X|^{2j}\right)^{1/2}.
\]
The uniform moment bounds and Lipschitz cumulant polynomials therefore
control each relevant power-sum discrepancy by $CH$. Apply
\cref{lem:moments}; for variance also use \eqref{eq:kappa2}.
The scaled pairs above have parameter gaps of order $h$ or $h^2$ and
Hellinger distances bounded above by order $h^{2m}$ or $h^{2m+2}$.
They exclude every larger exponent.
\end{proof}

\section{Detecting a factor is easier than recovering every factor}
\label{sec:detection}

Let $\mathcal T_n(h)$ be the infimum over tests of the sum of their
worst-case null and alternative errors, with null $a=0$ and alternative
$a_m\ge h$. In the unknown-variance model the supremum over $v\in[0,V]$
is taken separately under the two hypotheses. This is essential: the
null and alternative need not have the same nuisance variance.

\begin{proof}[Proof of \cref{thm:detection-main}]
With known $v_0$, define the estimated variance deficit
\[
 \widehat D_2=\widehat\kappa_2-\kappa_2(B)-v_0.
\]
Its true value is zero under the null and at least $h^2/3$ under the
alternative. Reject when $\widehat D_2>h^2/6$. The moment-estimation
bound \eqref{eq:kappa-MSE} and Markov's inequality applied to squared
error give
\begin{equation}
 \mathcal T_n(h)\le\frac{C}{nh^4}.
 \label{eq:test-upper-known}
\end{equation}
With unknown variance instead use
$\widehat D_4=\widehat\kappa_4-\kappa_4(B)$. Its true value is zero
under the null and at most $-2h^4/15$ under the alternative, because
$\kappa_4(X)-\kappa_4(B)=-(2/15)\sum_i a_i^4$.
Reject when $\widehat D_4<-h^4/15$. The same argument gives
\begin{equation}
 \mathcal T_n(h)\le\frac{C}{nh^8}.
 \label{eq:test-upper-unknown}
\end{equation}

For the known-variance lower bound, compare no factor with just one factor
of half-length $h$. At $v_0=0$, \eqref{eq:up-single} with $f$ in place
of $\up$ gives Hellinger squared asymptotic $J_2(f)h^4/144$.
For $v_0>0$, use common-convolution contraction. If $nh^4\to0$,
\eqref{eq:tensorization} implies that product total variation tends to
zero, forcing the error sum to be at least $1-o(1)$.

For the unknown-variance lower bound, use the null with $a=0$, $v=h^2/3$
and the alternative with one half-length $h$, $v=0$. Both lie in the
parameter set for small $h$. The perturbations have matching variance and
first differ in their fourth moment, so
\begin{equation}
 H^2(P_{(0,\ldots,0,h),0},P_{0,h^2/3})
 \sim\frac{J_4(f)}{129600}h^8.
 \label{eq:test-hard-pair}
\end{equation}
If $nh^8\to0$, product total variation again tends to zero. The trivial
always-accept test has error sum one, so in both lower-bound regimes the
optimal sum tends exactly to one. Together with
\eqref{eq:test-upper-known}--\eqref{eq:test-upper-unknown}, this proves
the threshold statements.
\end{proof}

For $m>1$, the worst-case recovery rate is strictly slower than the
corresponding detection threshold. The hard recovery pairs use several
coordinated small factors whose low cumulants cancel between the two
models. In a test against no factors, all contributions to the first
unconfounded even power sum have the same sign. There is no such
cancellation to exploit.

\section{Concrete two-factor hard pairs}
\label{sec:examples}

The implicit-function construction proves existence for every capacity.
For $m=2$ the following exact rational shapes make both mechanisms
transparent. These examples use squared half-lengths.

\subsection{Known variance}
Take
\begin{equation}
 u=(1/4,3/4),\qquad w=(1/8,7/8).
 \label{eq:example-known}
\end{equation}
Their first power sums agree and
$p_2(u)-p_2(w)=-5/32$. Hence their scaled uniform noises have first
unmatched raw moment coefficient
$\Delta_4=(-2/15)(-5/32)=1/48$. At zero Gaussian variance,
\begin{equation}
 H^2(P_{h\sqrt u,0},P_{h\sqrt w,0})
 \sim\frac{J_4(f)}{5308416}h^8.
 \label{eq:example-known-H}
\end{equation}
The two scale vectors differ by a fixed positive multiple of $h$, giving
the $n^{-1/8}$ two-factor recovery obstruction.

\subsection{Unknown variance, one exactly zero}
Keep $u=(1/4,3/4)$ and take
\begin{equation}
 w=(39/404,317/404).
 \label{eq:example-unknown}
\end{equation}
Direct calculation gives
\[
 p_2(u)=p_2(w)=5/8,\qquad p_1(u)=1,\qquad p_1(w)=89/101.
\]
Thus choose $t_u=0$, $t_w=4/101$. The second cumulants of the two
unscaled perturbations are both $1/3$, their fourth cumulants agree,
and their sixth raw-moment difference is
\begin{equation}
 \Delta_6=\frac{16}{63}\bigl(p_3(u)-p_3(w)\bigr)
 =-\frac{85164}{7212107}.
 \label{eq:example-delta6}
\end{equation}
Therefore
\begin{equation}
 H^2(P_{h\sqrt u,0},P_{h\sqrt w,(4/101)h^2})
 \sim\frac{\Delta_6^2}{4(6!)^2}J_6(f)h^{12}.
 \label{eq:example-unknown-H}
\end{equation}
This is a zero-noise-boundary realization of the $n^{-1/12}$ scale and
$n^{-1/6}$ variance obstructions for capacity two, not an appeal to a
positive variance floor.

\section{Computational checks and their limits}
\label{sec:verification}

The package includes \texttt{verification.py}, using SymPy 1.14.0 and
mpmath 1.3.0. The executed run passed 2,100 exact rational instances of
\eqref{eq:moment-inverse} for $1\le m\le7$ and $0\le s\le2$, using
seed 20260929; it also passed 16 exact tangent identities for
$1\le m\le8$, $s=0,1$. It verifies the first six uniform cumulants,
the displayed hard-pair equalities, and the rational Hellinger constants.
These checks independently test algebraic ingredients and transcription;
they do not replace the all-parameter proofs.

The numerical part uses 60-decimal arithmetic and piecewise integration at
every relevant spline knot. It computes squared Hellinger distances for
$g_M*\Law(hU)$ versus $g_M$, with $M=1,3,4,5$ and
$h=2^{-3},\ldots,2^{-7}$. There are twenty diagnostic rows. The finite
spline formula is \eqref{eq:spline}; convolution with the added uniform is
obtained by integrating each truncated power exactly before quadrature.
The full data are in \path{data/hellinger_diagnostics.csv}.

\begin{table}[htbp]
\centering
\small
\begin{tabular}{@{}rlll@{}}
\toprule
$M$ & Quantity at $h=1/128$ & Computed value & Predicted limit\\
\midrule
1 & $H^2/h$ & $0.276142374915$ & $0.276142374915$\\
3 & $H^2/h^3$ & $0.017668468430$ & $0.017661491088$\\
4 & logarithmic slope $L(h)$ & $0.005208430349$ & $1/192$\\
5 & $H^2/h^4$ & $0.007795167122$ & $0.007819937017$\\
\bottomrule
\end{tabular}
\caption{Finite-prefix diagnostics; numerical quadrature, not interval
certificates. Here $L(h)=[H^2(h)/h^4-H^2(2h)/(2h)^4]/\log2$.}
\label{tab:diagnostics}
\end{table}

For $M=3$, the predicted limit is $2C_3A_3=A_3/8$ and the independently
integrated profile is
\[
 A_3\approx0.141291928707536383107313127009.
\]
For $M=5$ the computation gives
\[
 J_2(g_5)\approx1.126070930490470658853603770783,
 \qquad J_2(g_5)/144\approx0.00781993701729493513.
\]
The critical $M=4$ ratio $H^2/(h^4\log(1/h))$ converges slowly because
of the next constant term. The slope $L(h)$ removes that leading constant
and is a more informative finite-$h$ diagnostic of the coefficient $1/192$.
The one-uniform formula \eqref{eq:M1} is also checked numerically against
its exact value at every tested amplitude.

No interval enclosure of an information integral is claimed. In particular,
none of the displayed finite-prefix numbers is described as a computed
$J_r(\up)$. The package does not simulate a minimax risk, fit the enormous
moment grid, or certify an asymptotic statement by finite sampling.
The exact asymptotic constants for the up law remain the rigorously finite,
positive integrals in the theorems.

\section{Proof audit and formalization route}
\label{sec:audit}

The dependency chain is short enough to audit structurally:
\[
 \begin{gathered}
 \text{finite spline endpoints and conditional Jensen}\\
 \Downarrow\\
 \text{finite information, prefix information convergence}\\
 \Downarrow\\
 \text{finite Hellinger expansion + contraction + Fatou}\\
 \Downarrow\\
 \text{exact infinite-background Hellinger asymptotic}\\
 \Downarrow\\
 \text{moment-matched pairs + tensorization + two-point risk}.
 \end{gathered}
\]
The statistical upper bounds use a separate branch: uniform moment bounds,
finite cumulant polynomials, and the credited inverse inequality.

Three distinctions are essential. First, $C^\infty$ smoothness alone is
not used as a theorem that a weighted likelihood expansion must hold.
The uniform-prefix representation supplies the missing quantitative control.
Second, passing to the infinite convolution uses convergence of information,
not just its lower semicontinuity; the higher-score moment bound in
\eqref{eq:uniform-integrability} is what establishes the required direction.
Third, no upper bound on chi-square divergence relative to the unsmoothed
law is asserted. It is actually infinite in a central example.

A useful proof-assistant implementation would separate four interfaces.
The first is the finite uniform-convolution density, support, and endpoint
formula. The second is weak differentiation of convolution, conditional
Jensen, and weighted uniform integrability. The third is Hellinger affinity,
its contraction, tensorization, and the prefix-limit argument. The fourth
is finite moment algebra and the statistical decision bound. A certified
checker for rational tangent identities could accompany the final interface,
but cannot replace the analytic interfaces.

No Lean or Rocq files are delivered as verified proofs. Nothing inherits
formal status merely from its association with the ProveIt repository.
Likewise, the article does not establish a full local asymptotic normality
theorem: the $L^2$ square-root expansion compares prescribed perturbation
curves, and identifying an entire limiting statistical experiment needs
additional uniformity.

\section{Further research questions and concrete success criteria}
\label{sec:questions}

\begin{question}[Mixed collision strata at the zero-noise boundary]
At a base factor configuration with positive squared-scale multiplicities
$m_1,\ldots,m_s$ and $m_0$ zero slots, determine the exact local pairwise
Hellinger inverse exponent with unknown variance. The companion manuscript
proposes the denominator
\[
 \max\bigl(\{m_1,\ldots,m_s\}\cup
 \{2(m_0+1):m_0>0\}\bigr).
\]
The present paper proves the fully vanishing case, not this entire formula.
A solution should supply both a local upper modulus and moment-matched
likelihood pairs for every dominant cluster, including $v=0$.
\end{question}
% ed. (2026-09-29): the positive-variance analogue is settled by a later
% package; the zero-noise question stays open.
\begin{ednote}
For Gaussian variance unknown in a compact interval bounded away from zero,
the later repository article \path{inverse-and-sampling/Local_Minimax_Geometry_Uniform_Factors/} (filed 2026-09-29, unreviewed;
its Corollary \texttt{cor:joint}) proves the proposed denominator for the
local minimax sampling rates. The zero-noise boundary asked for here,
including $v=0$, remains open; that article poses it as its own question on
mixed strata without a positive variance floor.
\end{ednote}

\begin{question}[Singular limiting experiments and optimal constants]
Describe the limiting experiments when
$a=n^{-1/(4m)}u$ with known zero variance, or
$(a,v)=(n^{-1/(4m+4)}u,n^{-1/(2m+2)}t)$ with unknown variance.
The exact Hellinger constants identify candidate information coefficients,
but do not alone identify the likelihood-ratio process or efficient risk
constant. A useful result would include convergence uniform over compact
sets of shape parameters and an attainable optimal constant.
\end{question}

\begin{question}[Certified up-information constants and their growth]
Develop interval-certified algorithms for $J_r(\up)$, beginning with
$r=2,4,6$. Monotone prefix values give upper bounds; interior integrals and
controlled endpoint estimates can give lower bounds. Prove an effective
error bound, then determine the growth of $J_r(\up)$ as $r\to\infty$.
Without such control, the fixed-capacity rates cannot be read uniformly in
an increasing number of factors.
\end{question}

\begin{question}[The finite-to-infinite crossover]
Let the available background prefix length $M=M(h)$ grow while $h\to0$,
or let the geometric ratio of $b_j$ vary with $h$. Determine when the
finite-boundary term, the critical logarithm, and the full information
constant dominate. A solution must track the width of the exact endpoint
collar and the size of the information constants; substituting a varying
$M$ into a fixed-$M$ asymptotic is not valid.
\end{question}

\begin{question}[Minimal tail hypotheses]
All absolute moments suffice in \cref{thm:main-hellinger}; sub-Gaussian
or exponential tails are unnecessary. For a prescribed order $r$, determine
the weakest finite-moment or tail condition ensuring the same expansion.
The truncation proof suggests a tradeoff between available prefix length
and available noise moments. Find sharp sufficient conditions and a
heavy-tail counterexample at the threshold.
\end{question}

\begin{question}[Beyond-all-powers separation]
For the explicit lognormal pair in \cref{cor:allmoments}, determine the
actual first nonzero asymptotic scale of the Hellinger distance. It is
smaller than every power of $h$, but nonzero. Does a logarithmic-square
scale analogous to the lognormal tail govern its leading behavior, and
how does the background enter the coefficient? This question lies beyond
finite-order moment methods by construction.
\end{question}

\begin{question}[Adaptation and honest uncertainty quantification]
Construct an estimator that adapts to unknown collision multiplicities,
vanishing slots, and positive versus zero variance. Determine whether
confidence sets can be simultaneously honest and rate-adaptive on all
these strata. Separate an impossibility caused by nearly indistinguishable
strata from inefficiency of a particular algorithm.
\end{question}

\begin{question}[Efficient certified moment fitting]
Replace the $\binom{n+m}{m}$ grid by an algorithm with explicit dependence
on $m$ and optimization tolerance. A successful certificate should guarantee
an objective gap $O(n^{-1/2})$, which is sufficient for the proved rates,
even when roots coincide or vanish. Hyperbolic-polynomial coefficient
methods and interval branch-and-bound are plausible starting points, not
implemented solutions in this package.
\end{question}

\begin{question}[Multivariate uniform directions]
Replace $a_iV_i$ by $V_i a_i$ with vectors $a_i\in\R^d$, and permit an
unknown Gaussian covariance matrix. Higher cumulants become sums of even
tensor powers. Determine identifiable parameter classes up to signs and
permutations, and then prove a multivariate finite-boundary phase diagram
that includes faces of different codimensions. The one-dimensional
counting-function inverse inequality does not extend automatically.
\end{question}

\begin{question}[Characterize admissible compact backgrounds]
Find a useful intrinsic condition on a compactly supported density $g$
that is equivalent, or close to equivalent, to the all-order Hellinger
moment-matching property. Is finite $J_r(g)$ together with an appropriate
translation-uniform integrability condition sufficient at each fixed $r$?
The present uniform-prefix representation is a strong, checkable sufficient
condition, but not claimed necessary.
\end{question}

\section{Conclusion}

For the unsmoothed Fabius--Rvachev law, a moving compact-support boundary
does not change the finite-order moment-based Hellinger exponent. The
reason is specific and quantitative: sufficiently long uniform prefixes
have enough endpoint vanishing to make the relevant score square
integrable, their information converges to the infinite-background
information, and common-convolution contraction transfers the sharp
asymptotic upper bound. Fatou supplies the matching lower bound.

This closes the identified unsmoothed statistical gap for every fixed
finite capacity and extends the lower-bound constructions to the boundary
of unknown Gaussian variance. Finite prefixes show why this conclusion
was not automatic: too few uniform factors produce a different power,
and the exact critical number produces a logarithm. The result is a
proved continuation of a specific research program, with its classical
algebra and information ingredients distinguished from the analytic and
zero-noise conclusions developed here.

\begin{thebibliography}{9}
\bibitem{arias}
J. Arias de Reyna,
\emph{An infinitely differentiable function with compact support:
Definition and properties}, arXiv:1702.05442 (2017), English translation
of the 1982 article in Rev. Real Acad. Ciencias Madrid \textbf{76}, 21--38.
\url{https://arxiv.org/abs/1702.05442}.

\bibitem{bobkov}
S. G. Bobkov,
\emph{Fisher-type information involving higher order derivatives},
arXiv:2412.10200v1 (2024), especially Sections 1, 2, and 10.
\url{https://arxiv.org/abs/2412.10200}.

\bibitem{heinrich-kahn}
P. Heinrich and J. Kahn,
\emph{Optimal rates for finite mixture estimation},
arXiv:1507.04313 (2015).
\url{https://arxiv.org/abs/1507.04313}.

\bibitem{wu-yang}
Y. Wu and P. Yang,
\emph{Optimal estimation of Gaussian mixtures via denoised method of moments},
arXiv:1807.07237v2 (2019).
\url{https://arxiv.org/abs/1807.07237v2}.

\bibitem{kleiber}
C. Kleiber,
\emph{The generalized lognormal distribution and the Stieltjes moment problem},
arXiv:1301.1277 (2013).
\url{https://arxiv.org/abs/1301.1277}.

\bibitem{repo-strata}
ProveIt research collection,
\emph{Sharp Stability Strata for Finite Fabius--Rvachev Deconvolution:
Colliding scales, vanishing factors, and optimal recovery exponents},
report dated 28 September 2026.
Read at commit \texttt{9b24a3a8d545af9624f6ac455f5b548be62818b6}.
Repository: \url{https://github.com/VladimirReshetnikov/ProveIt}.
Path: \path{Analysis/FabiusFunction/docs/semi-formalized-research-frontiers/drafts/inverse-and-sampling/Sharp_Stability_Strata_Fabius_Rvachev_Deconvolution/article.tex}.
The sampling proposition and explicit lower-bound limitation are in source
lines 1190--1234 at that snapshot. This is an unreviewed research report,
not a cited formal theorem.

\bibitem{companion}
\emph{Gaussian Confounding and Sharp Recovery of Uniform Convolution Factors:
Shifted power sums, minimax estimation, and detection above a
Fabius--Rvachev background}, research report prepared for Vladimir
Reshetnikov, 29 September 2026, 19 pages.
Read from the user's Library as \path{article(20260929-161307).pdf}, version 1.
Theorem 4.1 supplies the credited shifted-moment inequality; Research
Question 1 on page 17 explicitly asks for the unsmoothed up-law extension.
Its own repository comparison is pinned to
\texttt{c130dba623c420551d90db41dae7b3d9cff72dfd}, not the snapshot used
here. No publication or independent review is implied.
\end{thebibliography}
\end{document}
